\section{Introduction}

Large language models have recently demonstrated agentic capabilities---planning, tool invocation, and long-horizon reasoning---that make them attractive cores for intelligent assistants in electronic design automation (EDA).
% 中文：大语言模型近来展现出智能体能力——规划、工具调用与长程推理——这使其成为电子设计自动化（EDA）中智能助手的理想内核。
Such assistants are no longer hypothetical: agentic EDA flows are already deployed in-house, with a frontier large language model serving as the orchestrator that coordinates the flow.
% 中文：此类助手已非假设：智能体 EDA 流程已在公司内部署，由一个大模型担任 orchestrator 负责协调。
The binding constraint is therefore cost, dominated by two terms: every subtask routed to the orchestrator consumes tokens and serving budget that scale with flow volume, and because an agentic session accumulates context, serving cost is driven by the memory that grows with the session rather than by prompt throughput.
% 中文：因此当前约束是成本，由两项主导：每个路由到 编排者 的子任务都消耗随流程规模增长的 token 与 serving 预算；且由于智能体会话会累积上下文，服务成本由随会话增长的内存而非提示词吞吐驱动。
A small, locally served subagent caps the cost of every subtask it takes over by pricing it at a local 2B model rather than at an orchestrator.
% 中文：本地部署的小型子智能体把所接管的每个子任务按本地 22 亿参数模型而非 orchestrator 计价，从而为成本设上界。
This is the metaphor of the \emph{Artifact Spirit} (\qi): the orchestrator stays in the cloud to hold the user's intent and long-horizon plan, while every EDA tool is paired with its own small, local spirit that translates intent into tool-native scripts and operates the tool.
% 中文：这就是「器灵」（Artifact Spirit）的隐喻：编排者 留在云端把握用户意图与长程规划，而每一个 EDA 工具都配有自己的小型本地灵体，把意图翻译成工具原生脚本并操作该工具。
Privacy compounds the cost argument: netlists, GDSII layouts, and design rules are trade secrets that may not leave the premises, so for on-site subtasks the choice is between a small local model and no model at all.
% 中文：隐私进一步强化了成本论证：网表、GDSII 版图与设计规则是商业机密，不得离开本地，因此对本地子任务而言，选择只能是「小型本地模型」或「根本没有模型」。
Building such a model \emph{from scratch} as a production subagent, however, remains an under-documented engineering problem.

\paragraph{Contributions.}
\begin{itemize}
\item \textbf{A 2B Mamba2-hybrid backbone with a 1M-token context window.}
We build a $\approx$2B hybrid (state-space + attention) model from scratch whose 56-layer backbone interleaves Mamba-2 SSM and MLP layers with only four GQA attention layers (Section~\ref{sec:arch}), and whose RoPE configuration extrapolates to a $1$M-token context window with stable perplexity and ${\geq}95\%$ needle-retrieval accuracy (Section~\ref{sec:longctx}).
% 中文：2B Mamba2 混合骨干与 1M token 上下文窗口。我们从零构建一个约 22 亿参数的混合（状态空间+注意力）模型，其 56 层骨干将 Mamba-2 SSM 与 MLP 层同仅 4 个 GQA 注意力层交错排列（第 \ref{sec:arch} 节），并通过 RoPE 配置外推到 1M token 上下文窗口，困惑度稳定、针检索精度 $\geq95\%$（第 \ref{sec:longctx} 节）。

\item \textbf{Curriculum learning that strengthens code, math, and agentic capability.}
A three-stage curriculum---two stable phases (S1/S2; $700$k steps) and a decay annealing phase---shifts the data mixture from a general $74{:}19{:}7$ web/code/math blend to a code-and-math-enriched $40{:}40{:}20$ blend, then anneals on instruction and agentic SFT data with a long-context ladder (Section~\ref{sec:llm-pretrain}).
% 中文：通过课程学习强化代码、数学与智能体能力。三阶段课程学习——两个稳定阶段（S1/S2；700k 步）加一个退火阶段——将数据配比从通用 74:19:7（网页/代码/数学）逐步切换到强化代码与数学的 40:40:20 配比，再在指令与智能体 SFT 数据上退火并叠加长上下文阶梯（第 \ref{sec:llm-pretrain} 节）。

\item \textbf{Domain adaptation validating the subagent as a deployable EDA agent core.}
Through domain annealing, supervised fine-tuning, and reinforcement learning (Section~\ref{sec:llm-posttrain}), we adapt the pre-trained backbone into a subagent that serves EDA agents and assists circuit designers in orchestrating design flows---where EDA scripts serve as the glue between intent and tool execution---validating the feasibility of a large--small model collaborative agent architecture for agentic IC design.
% 中文：领域适配验证子智能体可作为可部署的 EDA Agent 内核。通过领域退火、监督微调与强化学习（第 \ref{sec:llm-posttrain} 节），我们把预训练骨干适配为子智能体，服务 EDA Agent、协助电路设计师编排 designflow（EDA 脚本充当意图与工具执行之间的胶水），验证「大小模型结合的 agent 架构在 agentic IC 设计落地的可行性」。
\end{itemize}
